% ===== §4 The Political Economy of Experience =====
\section{The Political Economy of Experience: Reproduction and the Circulation of Relational Wealth}
\label{p21:sec:reproduction}

\noindent
The field of travel generates something; the previous section described what. But generation is only half of any complete account, and on its own it is the half that dissipates. A value generated and not reproduced is a value squandered, present once and then gone, leaving no deposit. This section turns from the generation of relational experience to its \emph{reproduction}: the process by which what the field generates is made to persist, circulate, and accumulate rather than vanish with the journey's end. The turn is also the turn from phenomenology to political economy, for reproduction is a category of political economy before it is anything else, and the claim of this section is that the afterlife of a shared journey is best understood through it. The section argues that the experience a journey generates is a form of relational wealth; that its retelling is the labour that reproduces that wealth; that reproduction is the establishment of a good circulation; that this living labour is distinct from, and not to be confused with, the dead material (the photograph, the souvenir) that is merely its raw stuff; that the moment of return is the critical juncture at which reproduction either begins or fails; and that reproduction, like all reproduction, can be just or unjust in its distribution.

% ============================================================
\subsection{The experience of the journey is relational wealth}
\label{p21:sub:reproduction-wealth}

This series has already established the category of generative relational wealth: the value that a relation generates and holds not as a stock divided between two parties but as a genuinely common surplus, the wealth of the ``we'' rather than the summed possessions of the two ``I''s \parencite{paperxv}. The experience a shared journey generates is wealth of exactly this kind, and the point need not be argued at length because it follows directly: the self-presence of the relational dynamics, the communitas of the threshold, the depth reached in the cleared field, these are not goods that either partner possesses privately and could carry off alone, but goods of the relation, generated by and belonging to the ``we.'' What the two bring home from a journey that went well is a common deposit in the wealth of their relation rather than two private stores of pleasant memory. This is why the loss of it, when it is lost, is felt as a loss to the relation and not merely to two individuals; and it is why its conservation is a matter for the relation's political economy, the question of how the ``we'' manages the wealth it generates.

But relational wealth, like any wealth, is subject to a basic political-economic fact: it is not automatically conserved. Generated, it may be reproduced, in which case it persists and accumulates; or it may be left unreproduced, in which case it dissipates. The journey's experience does not deposit itself permanently into the relation's wealth merely by having occurred. It must be \emph{reproduced} into the relation's ongoing life, or it fades, and the couple who return from a journey that genuinely generated something, and then do nothing to reproduce it, will find within months that it has thinned to a few photographs they no longer look at and a fact they can state but no longer feel. Generation without reproduction is a one-time expenditure. The wealth is real, but it is perishable, and its conservation is a labour.

% ============================================================
\subsection{Retelling as the reproduction of subjective experience}
\label{p21:sub:reproduction-retelling}

What is that labour? It is, principally, \emph{retelling}: the shared, spoken, mutual re-narration of the journey, by which the two return to what they generated and weave it, again and again, into the fabric of the ``we.'' \emph{Do you remember that time we got lost}, the small phrase is the beginning of a great deal, for in the shared retelling the generated experience is not merely recalled but \emph{re-produced}: brought back into the relation's living present, re-felt, re-woven, and thereby kept from dissipating. This is the \emph{reproduction of subjective experience}, and the name is chosen with its full political-economic weight. In the political economy of labour, reproduction is the process by which what production has generated is made to persist, circulate, and serve as the precondition of the next round of production; it is the condition of production's being more than a single rather than an appendage to production, dissipating event \parencite{federici2004}. Reproduction, here, is the same category applied to a different object: not the reproduction of labour-power but the reproduction of the relational subject's experience, the labour by which a generated experience is kept in circulation and thereby conserved as wealth.

The application is exact in each of its parts. Production, the generation of the experience in the field, is, on its own, a one-time event; the journey ends, the field closes, and what was generated would dissipate were it not reproduced. Reproduction, the shared retelling, is what keeps it: by returning the experience to the relation's present, the retelling makes it persist beyond its moment, circulate through the relation's ongoing life, and become available as the precondition of further generation, the deposited wealth on which the relation can draw. And reproduction is, like all reproduction, \emph{labour}, living activity, expended by the two, without which the conservation does not happen. This last point is the one most easily missed and most important, and the next subsection is devoted to it; for the great temptation of the age is to imagine that the reproduction of experience has been automated, outsourced to a device, when in truth it cannot be performed by anything but the living labour of the two who shared the experience.

% ============================================================
\subsection{Reproduction is the establishment of a good cycle}
\label{p21:sub:reproduction-cycle}

The retelling does not merely preserve; it \emph{circulates}, and in circulating it can do something that mere preservation cannot, it can establish a good cycle, in the exact sense this series has given the term. A good cycle, as developed in the foundational treatment of sustainable intimacy, is a circulation that returns more than it took: a loop along which positive surplus accumulates, so that the relation comes back to itself not depleted but enriched, the circulation generating rather than merely conserving \parencite{paperix}. The shared retelling, when it is alive, is such a cycle. Each return to the journey in retelling is not a static replay of a fixed content but a re-generation: the experience is re-felt, and in the re-feeling it is added to, inflected by all that has happened since, enriched by the very act of being retold, so that the couple who return often to a shared journey do not wear it out but deepen it, the story growing in meaning each time it is told. The retelling is a loop that returns more than it took, the holonomy of the relation accumulating positive phase along the path of re-narration, to put it in the terms of the series' formal companion, where the surplus of a relation is exactly what a closed circulation generates over its inputs \parencite{wan2026braid}. Reproduction, then, is the establishment of a circulation rather than the mere maintenance of a stock, and a circulation that, when healthy, is generative: the good cycle by which a journey's wealth not only persists but grows.

This is the deep reason the reproduction of experience matters as much as its generation, and perhaps more. Generation gives the relation a wealth; reproduction establishes the circulation by which that wealth becomes self-augmenting, by which a single journey, retold and re-felt over years, becomes an ever-deepening common possession rather than a fading memory. A relation rich in such circulations, many shared experiences, each kept in living retelling, is rich in the most durable way, for its wealth is a set of good cycles each of which returns more than it took rather than a fixed store that can only be spent down. The political economy of the relation is, in large part, the management of these cycles: their establishment, their maintenance, and the subject of the next subsections, their two characteristic failures, reification and unjust distribution.

% ============================================================
\subsection{Dead material and living labour: the photograph}
\label{p21:sub:reproduction-deadlabour}

The reproduction of experience is living labour, and it cannot be performed by anything else. This must be insisted upon, because the most pervasive error in the modern conduct of travel is the belief that the reproduction has been accomplished by the \emph{recording}, that to have photographed the journey is to have preserved it. The photograph, the video, the souvenir, the geotagged log: these are not the reproduction of the experience. They are, at most, its \emph{material}, the dead stuff from which a reproduction might later be made, but which is not itself a reproduction and accomplishes nothing on its own. The distinction is precisely Marx's distinction between dead and living labour. Dead labour is labour congealed in a thing, past activity solidified into an object; living labour is present activity, and only living labour can set dead labour back into motion and make it yield \parencite{marx1976}. The photograph is dead labour: the congealed residue of a past seeing, a thing in which an experience has been deposited and frozen. It can become material for reproduction, a prompt, a prop, a trigger for the living retelling, but only if living labour takes it up and works it. A photograph never returned to in shared retelling reproduces nothing; it is dead labour that stays dead, a stock of congealed seeing accumulating on a drive, untouched by the living activity that alone could make it yield.

The consequences of confusing the material with the reproduction are severe, and they are visible everywhere in the contemporary conduct of travel. The couple who photograph everything and retell nothing have mistaken the accumulation of dead material for the conservation of wealth, and they will find their wealth dissipating beneath a growing pile of unviewed images, the perfect emblem of dead labour accumulating while the living labour that alone could redeem it is never performed. Worse, the labour of recording can actively \emph{displace} the labour of reproduction: the energy that goes into capturing the moment for later is energy not spent undergoing the moment now or retelling it after, and the camera interposed between the two and their experience can prevent the very generation it was meant to preserve. The journey lived through a lens, photographed rather than undergone, may generate little to reproduce; and the journey whose afterlife is a folder of images never reopened reproduces little of what it generated. The recording is not the reproduction. Only the living labour of shared retelling reproduces the experience, and the dead material is worth exactly as much as the living labour that returns to work it, and not one bit more.

% ============================================================
\subsection{The return: the critical juncture}
\label{p21:sub:reproduction-return}

There is a moment at which reproduction either begins or fails, and it has been underweighted in every account of travel that treats the journey as ending at the homecoming. That moment is the \emph{return} itself, the descent from the field back into the everyday, the crossing of the threshold in the homeward direction. The return is not the journey's conclusion; it is the journey's most delicate phase, the juncture at which everything the field generated must be carried across into a structure that is hostile to it, or be left behind at the threshold. The hostility is structural, established already in the analysis of the field's tensions (\xref{p21:sub:field-tensions}): the communitas of the journey is a liminal phenomenon, anti-structural by nature, and the everyday into which one returns is precisely structure, the reimposition of the roles that suspension had bracketed, the predetermination that contingency had lifted, the channelled time that reshaped time had opened. To return is to step from the threshold back into the grid, and the grid does not preserve what the threshold generated; left to itself, it dissolves it. This is the structural truth behind the most common lament of travellers, \emph{it was so good there, and within a week of being home we were back to how we always are}. The lament is not a failure of feeling but a failure of reproduction at the critical juncture: the field's wealth was generated and then surrendered at the threshold, because nothing was done to carry it across.

The return therefore is its beginning rather than the end of the labour. To carry the field's wealth across the threshold is to begin, immediately and deliberately, the reproduction that will let it survive in the structured everyday: to retell while the experience is fresh, to deposit it into the relation's living present before the grid can dissolve it, to begin the good cycle while the wealth is still warm. And, this is the hinge on which the paper's final movement will turn, the labour of carrying the field's wealth across the threshold into the everyday is continuous with the labour of constructing the field within the everyday. To reproduce a journey's wealth at home is already to begin building, in the everyday, a small continuation of the field that generated it; the return is where the literal journey begins to become the generalised one (\xref{p21:sec:generalized}). The couple who reproduce well do not simply remember their journey; they let it seed the construction of fields in the everyday, so that the homecoming is the carrying of the field's capacity home rather than a fall from the field. The return is the juncture where this either happens or does not, where the journey either deposits its wealth and seeds its own continuation, or is surrendered at the threshold and mourned.

% ============================================================
\subsection{The justice of reproduction}
\label{p21:sub:reproduction-justice}

Reproduction, finally, has a politics, because the retelling that reproduces the experience is never neutral, and the question of \emph{whose} version becomes the reproduced experience is a question of justice. Two people do not undergo the same journey identically: what one felt as a romantic adventure of getting lost, the other may have felt as a frightening consequence of the first's carelessness; what one remembers as a glorious meal, the other remembers as the evening they were ignored. The shared retelling weaves these divergent undergoings into a single shared narrative, and the weaving has a distribution of power. If one partner's version systematically dominates the retelling, if the official memory of the ``we'' is consistently the experience of one and not the other, then the reproduced experience is the product of an \emph{unequal reproduction}, a narrative appropriation in which one party's subjective experience is taken up as the relation's memory while the other's is overwritten or silenced. This is, in the register of memory, an exact instance of the bad cycle this series has analysed throughout: one party rendered the mere material of a circulation whose surplus accrues to the other, fuel for a wealth not jointly held \parencite{paperix,paperxvii}. The unjustly reproduced journey enriches the ``we'' in name while in fact enriching one of the two at the other's expense, the silenced partner's experience consumed to feed a common memory that is not truly common.

The just form of reproduction is, correspondingly, the one this series has described under the name of mutual translation: the retelling in which both undergoings are admitted, in which the two versions move toward each other and generate a shared narrative that neither alone possessed, the memory of the ``we'' built from both rather than imposed by one \parencite{paperxiii}. Just reproduction is harder than unjust, as mutual translation is harder than imposition; it requires that each partner's experience be given standing in the common memory, that the divergences be acknowledged rather than smoothed over by the dominant version, that the retelling generate a genuinely shared narrative rather than ratify one party's. But only just reproduction conserves the wealth as \emph{relational} wealth, the wealth of the ``we''; unjust reproduction converts it into the private wealth of one, dressed in the first-person plural. The justice of reproduction is thus not a separate ethical overlay upon a politically-economic process but internal to the process itself: only the just retelling reproduces the relation's wealth as the relation's, and the unjust retelling, whatever it preserves, has already begun to dissolve the ``we'' whose wealth it claims to keep. It is to the ethics implicit in this and the other operations of the field that the argument now turns.
